\chapter{चुंबकत्व और चुंबकीय क्षेत्र}\label{ch:g11:magnetic-fields}

नाविक हज़ार साल तक चुंबकित सुइयों के सहारे रास्ता खोजते रहे, इससे पहले कि
कोई जान पाता कि वे उत्तर की ओर क्यों मुड़ जाती हैं। उत्तर 1820 में मिला, जब
एक धारा ने कुतुबनुमा हिला दिया: चुंबकत्व चलते आवेशों से बनता है। यह अध्याय
\omterm{def:g11:magnetic-fields:bfield}{चुंबकीय क्षेत्र} का नक्शा बनाता है --- चुंबकों का, पृथ्वी का,
धाराओं का --- और फिर फंदा बंद करता है: क्षेत्र धाराओं पर पलटकर
धकेलता है, और हर विद्युत मोटर इसी से घूमती है।

\section{चुंबक और चुंबकीय क्षेत्र}

\begin{definition}[चुंबक और ध्रुव]\label{def:g11:magnetic-fields:magnet}
\emph{चुंबक}\index{चुंबक} लोहे को खींचता है और धागे से लटकाने पर घूमकर उत्तर
की ओर मुँह कर लेता है। उसकी क्रिया दो \emph{चुंबकीय ध्रुवों}\index{चुंबकीय
ध्रुव} पर सिमट जाती है, \emph{उत्तर} (वह सिरा जो भौगोलिक उत्तर खोजता है) और
\emph{दक्षिण}: समान ध्रुव प्रतिकर्षित करते हैं, विपरीत आकर्षित। चुंबक को बीच
से चीर दीजिए तो दो पूरे चुंबक मिलते हैं --- किसी प्रयोग ने आज तक अकेला ध्रुव
अलग नहीं किया।
\end{definition}

\begin{definition}[चुंबकीय क्षेत्र]\label{def:g11:magnetic-fields:bfield}
\omterm{def:g11:magnetic-fields:magnet}{चुंबक} अपने चारों ओर की जगह बदल देता है (\cref{ch:g11:electric-gravitational-fields}): वह हर जगह एक
\emph{चुंबकीय क्षेत्र}\index{चुंबकीय क्षेत्र} $\vect B$ रच देता है। कुतुबनुमा
उसकी दिशा पढ़ता है (सुई $\vect B$ के अनुदिश बैठ जाती है, उत्तरी सिरा आगे); और
\emph{हॉल प्रोब}\index{हॉल प्रोब} उसका परिमाण \emph{टेस्ला}\index{टेस्ला}
(\unit{T}) में मापता है।
\end{definition}

\begin{definition}[चुंबकीय क्षेत्र रेखाएँ]\label{def:g11:magnetic-fields:lines}
\emph{चुंबकीय क्षेत्र रेखाएँ}\index{चुंबकीय क्षेत्र रेखा} वे वक्र हैं जो
$\vect B$ के स्पर्शी हैं और जहाँ क्षेत्र प्रबल हो वहाँ सट जाती हैं।
\omterm{def:g11:magnetic-fields:magnet}{चुंबक} के बाहर वे उत्तरी ध्रुव से दक्षिणी ध्रुव तक जाती हैं;
\omterm{def:g11:electric-gravitational-fields:efield}{विद्युत क्षेत्र} की रेखाओं से उलट इनका कोई सिरा नहीं होता: हर रेखा
\omterm{def:g11:magnetic-fields:magnet}{चुंबक} के भीतर से होकर बंद हो जाती है।
\end{definition}

\begin{omfigure}
\begin{tikzpicture}[scale=0.9,
  fline/.style={omThm, postaction={decorate},
    decoration={markings, mark=at position 0.5 with {\arrow{Stealth}}}}]
  \fill[omProp!15] (-1.4,-0.35) rectangle (0,0.35);
  \fill[omDef!15] (0,-0.35) rectangle (1.4,0.35);
  \draw[thick] (-1.4,-0.35) rectangle (1.4,0.35);
  \draw (0,-0.35) -- (0,0.35);
  \node[font=\small] at (-0.7,0) {S};
  \node[font=\small] at (0.7,0) {N};
  \draw[fline] (1.4,0.2) .. controls (2.6,0.9) and (-2.6,0.9) .. (-1.4,0.2);
  \draw[fline] (1.4,0.3) .. controls (3.8,2.0) and (-3.8,2.0) .. (-1.4,0.3);
  \draw[fline] (1.4,-0.2) .. controls (2.6,-0.9) and (-2.6,-0.9) .. (-1.4,-0.2);
  \draw[fline] (1.4,-0.3) .. controls (3.8,-2.0) and (-3.8,-2.0) .. (-1.4,-0.3);
  \draw[fline] (1.4,0) -- (2.95,0);
  \draw[omExo, thick] (0,1.1) circle (0.22);
  \draw[-Stealth, omDef, thick] (0.22,1.1) -- (-0.22,1.1);
  \draw[omExo, thick] (2.45,0) circle (0.22);
  \draw[-Stealth, omDef, thick] (2.23,0) -- (2.67,0);
\end{tikzpicture}

{\small छड़ \omterm{def:g11:magnetic-fields:magnet}{चुंबक} की \omterm{def:g11:electric-gravitational-fields:lines}{क्षेत्र रेखाएँ}: N से बाहर, S में भीतर,
और शरीर के भीतर से बंद। कहीं भी रखा कुतुबनुमा उसी रेखा के अनुदिश ठहर जाता
है जो उससे होकर गुज़रती है।}
\end{omfigure}

\section{पृथ्वी एक चुंबक है}

\begin{definition}[पृथ्वी का चुंबकीय क्षेत्र]\label{def:g11:magnetic-fields:earth}
पृथ्वी के पिघले लोहे के \omterm{def:g10:refraction:fiber}{क्रोड} में बहती धाराएँ इस ग्रह को एक विशाल
\omterm{def:g11:magnetic-fields:magnet}{चुंबक} बना देती हैं: सतह पर $B \approx \qty{50}{\micro T}$, और स्वतंत्र सुई क्षेत्र
के क्षैतिज घटक के अनुदिश बैठ जाती है --- यही कुतुबनुमा है। भूचुंबकीय ध्रुव
भौगोलिक ध्रुव नहीं हैं; कुतुबनुमा के उत्तर और सच्चे उत्तर के बीच के कोण को
\emph{दिक्पात}\index{दिक्पात} कहते हैं, और हर दिशा-निर्धारण में उसका सुधार
किया जाता है।
\end{definition}

\begin{example}[दिशा का सुधार]\label{ex:g11:magnetic-fields:bearing}
जहाँ \omterm{def:g11:magnetic-fields:earth}{दिक्पात} $\ang{7}$ पश्चिम है, वहाँ ``कुतुबनुमा का उत्तर'' पकड़कर
\qty{5.0}{km} चलता जहाज़ अपने रास्ते से $5000 \times \sin\ang{7} \approx
\qty{6.1e2}{m}$ पश्चिम बह जाता है: सच्चे उत्तर की ओर
चलने के लिए सुई से $\ang{7}$ \emph{पूरब} हटकर चलिए।
\end{example}

\section{धारा का चुंबकीय क्षेत्र}

\begin{remark}[ओर्स्टेड, 1820]\label{rem:g11:magnetic-fields:oersted}
सन 1820 में व्याख्यान देते समय हैंस क्रिश्चियन ओर्स्टेड ने देखा कि पास ही
परिपथ बंद करते ही कुतुबनुमा की सुई घूम गई: बिजली और चुंबकत्व एक ही विषय निकले
--- \emph{धाराएँ \omterm{def:g11:magnetic-fields:bfield}{चुंबकीय क्षेत्र} रचती हैं}।
\end{remark}

\begin{proposition}[सीधे तार का क्षेत्र]\label{prop:g11:magnetic-fields:wire}
धारा $I$ ले जाता लंबा सीधा तार ऐसा क्षेत्र रचता है जिसकी रेखाएँ
तार पर केंद्रित वृत्त हैं और तार के लंबवत तलों में पड़ी हैं। तार से दूरी
$d$ पर,
\[
B = \frac{\mu_0\, I}{2\pi d},
\qquad \mu_0 = 4\pi\times 10^{-7}\ \unit{T.m/A} \approx \qty{1.26e-6}{T.m/A}.
\]
\end{proposition}
\admitted

\begin{omfigure}
\begin{tikzpicture}[scale=0.85]
  \node[font=\large] at (0,0) {$\odot$};
  \foreach \r in {0.55,1.05,1.55}
    \draw[omThm,-Stealth] (\r,0) arc (0:350:\r);
  \node[omThm, font=\small] at (45:1.9) {$\vect B$};
  \node[font=\small] at (0,-2.1) {धारा आपकी ओर};
\end{tikzpicture}\qquad\qquad
\begin{tikzpicture}[scale=0.85]
  \node[font=\large] at (0,0) {$\otimes$};
  \foreach \r in {0.55,1.05,1.55}
    \draw[omThm,-Stealth] (\r,0) arc (360:10:\r);
  \node[omThm, font=\small] at (45:1.9) {$\vect B$};
  \node[font=\small] at (0,-2.1) {धारा आपसे दूर};
\end{tikzpicture}

{\small सीधे तार के चारों ओर वृत्ताकार \omterm{def:g11:electric-gravitational-fields:lines}{क्षेत्र रेखाएँ}, सामने से देखी
गईं ($\odot$: धारा पन्ने से बाहर; $\otimes$: पन्ने के भीतर)। दाएँ हाथ से तार को
पकड़िए, अँगूठा धारा के अनुदिश: उँगलियाँ उसी ओर मुड़ती हैं जिधर $\vect B$ घूमता
है।}
\end{omfigure}

\begin{example}[तार एक कमज़ोर चुंबक है]\label{ex:g11:magnetic-fields:wire}
$I = \qty{15}{A}$ ले जाते तार से $d = \qty{1.0}{cm}$ पर: $B = \num{2e-7} \times 15 / 0.010 = \qty{0.30}{mT}$ --- यानी पृथ्वी के छह
क्षेत्र, और वह भी उस धारा से जो घर का ब्रेकर गिरा देगी। प्रबल
क्षेत्र के लिए बेहतर ज्यामिति चाहिए।
\end{example}

\begin{definition}[कुंडली और परिनालिका]\label{def:g11:magnetic-fields:solenoid}
तार को लपेट देने से उसका क्षेत्र सिमट जाता है। \emph{चपटी
कुंडली}\index{कुंडली} (एक चकती में $N$ फेरे) \omterm{def:g11:lenses-and-eye:lens}{पतले} \omterm{def:g11:magnetic-fields:magnet}{चुंबक}
जैसी बरतती है, एक फलक उत्तर और दूसरा दक्षिण; और
\emph{परिनालिका}\index{परिनालिका} --- लंबी बेलनाकार लपेट, प्रति \omterm{def:g10:orders-of-magnitude:unit}{मीटर}
$n$ फेरे --- छड़ \omterm{def:g11:magnetic-fields:magnet}{चुंबक} जैसी।
\end{definition}

\begin{proposition}[परिनालिका के भीतर का क्षेत्र]\label{prop:g11:magnetic-fields:solenoidfield}
लंबी \omterm{def:g11:magnetic-fields:solenoid}{परिनालिका} के भीतर, सिरों से दूर, क्षेत्र
\emph{\omterm{def:g11:electric-gravitational-fields:capacitor}{एकसमान}} रहता है: अक्ष के समांतर और भीतर के हर बिंदु पर एक-सा, जिसका
परिमाण $B = \mu_0\, n\, I$ है और जो नली की चौड़ाई पर निर्भर नहीं करता। बाहर
क्षेत्र दुर्बल है।
\end{proposition}
\admitted

\begin{remark}
दोनों सूत्र वर्ष 1 के खंड में उस नियम से ईमानदारी से निकाले गए हैं जो
क्षेत्र के परिसंचरण को उस धारा से जोड़ता है जिसे वह घेरता है। जो बात
मायने रखती है वह यह: $n$, यानी \emph{प्रति \omterm{def:g10:orders-of-magnitude:unit}{मीटर}} फेरे --- और
कसकर लपेटिए, लंबा मत कीजिए।
\end{remark}

\begin{method}[दायाँ हाथ, दो बार]\label{met:g11:magnetic-fields:righthand}
\begin{itemize}
  \item \emph{तार}: अँगूठा = धारा; मुड़ी उँगलियाँ = रेखाओं का घूमना।
  \item \emph{\omterm{def:g11:magnetic-fields:solenoid}{कुंडली}, \omterm{def:g11:magnetic-fields:solenoid}{परिनालिका}}: मुड़ी उँगलियाँ = फेरों में
        धारा; अँगूठा = भीतर $\vect B$, जो उत्तरी फलक की ओर संकेत करता है।
\end{itemize}
\end{method}

\begin{omfigure}
\begin{tikzpicture}[scale=0.9]
  \foreach \x in {0,0.65,...,3.9} {
    \node[font=\small] at (\x,0.85) {$\odot$};
    \node[font=\small] at (\x,-0.85) {$\otimes$};}
  \draw[-Stealth, omThm, thick] (-0.4,0.4) -- (4.3,0.4);
  \draw[-Stealth, omThm, thick] (-0.4,0) -- (4.3,0)
    node[right, font=\small] {$\vect B$};
  \draw[-Stealth, omThm, thick] (-0.4,-0.4) -- (4.3,-0.4);
  \draw[omExo, -Stealth, dashed]
    (4.5,0.7) .. controls (5.9,2.3) and (-2.0,2.3) .. (-0.6,0.7);
\end{tikzpicture}

{\small लंबाई के अनुदिश काटी गई \omterm{def:g11:magnetic-fields:solenoid}{परिनालिका} ($\odot$/$\otimes$: पन्ने को
काटते फेरे): भीतर अक्ष के अनुदिश \omterm{def:g11:electric-gravitational-fields:capacitor}{एकसमान} $B = \mu_0 n I$; बाहर दुर्बल, फैला हुआ लौटता
क्षेत्र (बिंदुदार)।}
\end{omfigure}

\begin{example}[परिनालिका के आँकड़े]\label{ex:g11:magnetic-fields:solenoid}
प्रति सेंटीमीटर $10$ फेरों ($n = \qty{1000}{m^{-1}}$) और $I = \qty{5.0}{A}$ पर: $B = \num{1.26e-6} \times 1000 \times 5.0 \approx
\qty{6.3}{mT}$ --- यानी पृथ्वी
के सौ क्षेत्र, और वह भी एक घुंडी से घटाया-बढ़ाया जा सकने वाला।
\end{example}

\section{विद्युत्चुंबक}

\begin{definition}[विद्युत्चुंबक]\label{def:g11:magnetic-fields:electromagnet}
\emph{विद्युत्चुंबक}\index{विद्युत्चुंबक} नरम लोहे के \omterm{def:g10:refraction:fiber}{क्रोड} पर लपेटी
\omterm{def:g11:magnetic-fields:solenoid}{परिनालिका} है: लोहा \omterm{def:g11:magnetic-fields:solenoid}{कुंडली} के क्षेत्र में चुंबकित हो
जाता है और उसे कई गुना कर देता है, आम तौर पर $100$ से $1000$ गुना। स्थायी
\omterm{def:g11:magnetic-fields:magnet}{चुंबक} से उलट, धारा बंद करते ही वह भी बंद हो जाता है।
\end{definition}

\begin{remark}[तीन मशीनें]\label{rem:g11:magnetic-fields:machines}
कबाड़खाने की क्रेन \omterm{def:g11:magnetic-fields:electromagnet}{विद्युत्चुंबक} से कार उठाती है और --- असली बात यही
है --- स्विच खोलकर उसे \emph{गिरा} देती है। रिले का छोटा
\omterm{def:g11:magnetic-fields:electromagnet}{विद्युत्चुंबक} एक संपर्क को खींचकर बंद कर देता है: यानी कमज़ोर धारा
प्रबल धारा को चलाती-बंद करती है (\cref{ch:g11:circuits-and-power})। और MRI यंत्र अतिचालक
\omterm{def:g11:magnetic-fields:solenoid}{परिनालिका} से रोगी के चारों ओर \omterm{def:g11:electric-gravitational-fields:capacitor}{एकसमान} \qty{1.5}{T} बनाए रखता है।
\end{remark}

आकाशगंगा से लेकर मरे हुए तारे तक, परिमाण की कोटियाँ:
\begin{center}
\small
\begin{tabular}{lr}
अंतरतारकीय अवकाश & $\sim \qty{e-10}{T}$ \\
पृथ्वी की सतह का क्षेत्र & \qty{5e-5}{T} \\
फ़्रिज का \omterm{def:g11:magnetic-fields:magnet}{चुंबक} & $\sim \qty{5e-3}{T}$ \\
नियोडिमियम \omterm{def:g11:magnetic-fields:magnet}{चुंबक}, उसके ध्रुव पर & $\sim \qty{0.5}{T}$ \\
MRI की \omterm{def:g11:magnetic-fields:solenoid}{परिनालिका} & \qty{1.5}{T} \\
न्यूट्रॉन तारे की सतह & $\sim \qty{e8}{T}$ \\
\end{tabular}
\end{center}

\section{लाप्लास बल}

\begin{proposition}[लाप्लास बल]\label{prop:g11:magnetic-fields:laplace}
लंबाई $L$ का सीधा तार, जो तार के \emph{लंबवत} \omterm{def:g11:electric-gravitational-fields:capacitor}{एकसमान क्षेत्र}
$\vect B$ में धारा $I$ ले जा रहा हो, $F = I\,L\,B$ परिमाण का \emph{लाप्लास
बल}\index{लाप्लास बल} महसूस करता है, जो तार और $\vect B$ दोनों के लंबवत है:
दायाँ हाथ चपटा रखिए, अँगूठा धारा के अनुदिश, सीधी उँगलियाँ $\vect B$ के अनुदिश
--- हथेली $\vect F$ के अनुदिश धकेलती है। क्षेत्र के \emph{समांतर} तार
पर कोई बल नहीं लगता।
\end{proposition}
\admitted

\begin{remark}[उछलती पटरी]\label{rem:g11:magnetic-fields:rail}
वर्ष 1 का खंड इसे हर चलते आवेश पर लगते चुंबकीय बल से निकालता है।
प्रयोगशाला में: ताँबे की एक छड़ किसी \omterm{def:g11:magnetic-fields:magnet}{चुंबक} के ध्रुवों के बीच दो
क्षैतिज पटरियों पर आड़ी रखी है (चित्र में क्षेत्र पन्ने के भीतर की
ओर); स्विच बंद कीजिए, और छड़ पटरियों पर सरपट भाग जाती है। धारा उलट दीजिए, या
\omterm{def:g11:magnetic-fields:magnet}{चुंबक} पलट दीजिए: वह दूसरी ओर भागती है।
\end{remark}

\begin{omfigure}
\begin{tikzpicture}[scale=0.95]
  \foreach \x in {0.6,1.5,2.9,3.8}
    \foreach \y in {0.35,1.05}
      \node[omThm, font=\small] at (\x,\y) {$\otimes$};
  \node[omThm, font=\small] at (4.35,1.05) {$\vect B$};
  \draw[very thick] (0,0) -- (4.6,0);
  \draw[very thick] (0,1.4) -- (4.6,1.4);
  \draw[omDef, line width=2.5pt] (2.1,-0.12) -- (2.1,1.52);
  \draw[-Stealth, omDef, thick] (0.1,1.6) -- (1.6,1.6)
    node[midway, above, font=\small] {$I$};
  \draw[-Stealth, omDef, thick] (2.32,1.15) -- (2.32,0.45);
  \draw[-Stealth, omDef, thick] (1.6,-0.22) -- (0.1,-0.22);
  \draw[-Stealth, omMeth, very thick] (2.1,0.7) -- (3.4,0.7)
    node[pos=0.85, above, font=\small] {$\vect F$};
\end{tikzpicture}

{\small पटरी वाला प्रयोग: धारा छड़ में नीचे की ओर, क्षेत्र पन्ने के
भीतर ($\otimes$) --- \omterm{prop:g11:magnetic-fields:laplace}{लाप्लास बल} $F = ILB$ छड़ को पटरियों पर दौड़ा देता
है।}
\end{omfigure}

\begin{example}[लाप्लास के आँकड़े]\label{ex:g11:magnetic-fields:laplace}
$B = \qty{0.50}{T}$ के क्षेत्र में $I = \qty{8.0}{A}$ ले जाती $L = \qty{5.0}{cm}$ लंबी छड़: $F = 8.0 \times 0.050 \times 0.50 = \qty{0.20}{N}$ --- यानी \qty{20}{g}
का भार, और वह भी कुछ ग्राम भारी तार पर।
\end{example}

\begin{remark}[दिष्ट धारा मोटर]\label{rem:g11:magnetic-fields:motor}
तार को मोड़कर क्षेत्र में एक फंदा बना दीजिए: $\vect B$ के लंबवत उसकी
दोनों भुजाओं में धाराएँ विपरीत हैं, इसलिए उनके \omterm{prop:g11:magnetic-fields:laplace}{लाप्लास बल} भी विपरीत
--- और यह युग्म फंदे को घुमा देता है। आधा चक्कर बाद वही युग्म उसे उलटा घुमाने
लगेगा, इसलिए एक घूमता स्विच (\emph{दिक्परिवर्तक}\index{दिक्परिवर्तक}) हर आधे
चक्कर पर धारा उलट देता है। हर पंखा और हर ड्रिल यही फंदा है, बस कई गुना।
\end{remark}

\section{अभ्यास}

\begin{exercise}[$\star$]\label{exo:g11:magnetic-fields:1}
छड़ \omterm{def:g11:magnetic-fields:magnet}{चुंबक} को उसके ध्रुवों के बीच से चीर दिया जाता है। हर टुकड़े पर
कौन-से ध्रुव हैं? फिर से पास लाने पर ताज़े कटे फलक आकर्षित करते हैं या
प्रतिकर्षित? क्या काटने का कोई तरीक़ा उत्तरी ध्रुव को अलग कर सकता है?
\end{exercise}

\begin{exercise}[$\star$]\label{exo:g11:magnetic-fields:2}
सही या ग़लत, हर एक के लिए एक कारण के साथ: (क) जहाँ क्षेत्र प्रबल
हो, वहाँ दो \omterm{def:g11:magnetic-fields:lines}{चुंबकीय क्षेत्र रेखाएँ} कट जाती हैं; (ख) \omterm{def:g11:magnetic-fields:magnet}{चुंबक}
के बाहर \omterm{def:g11:electric-gravitational-fields:lines}{क्षेत्र रेखाएँ} उत्तरी ध्रुव से दक्षिणी ध्रुव तक जाती हैं;
(ग) हर \omterm{def:g11:magnetic-fields:lines}{चुंबकीय क्षेत्र रेखा} बंद फंदा होती है; (घ) सुई अपने से होकर
जाती रेखा के लंबवत ठहरती है।
\end{exercise}

\begin{exercise}[$\star$]\label{exo:g11:magnetic-fields:3}
\qty{10}{A} ले जाते सीधे तार से \qty{2.0}{cm} पर क्षेत्र निकालिए। यह पृथ्वी
के \qty{50}{\micro T} का कितना गुना है?
\end{exercise}

\begin{exercise}[$\star$]\label{exo:g11:magnetic-fields:4}
\qty{40}{cm} पर लपेटे $800$ फेरों की \omterm{def:g11:magnetic-fields:solenoid}{परिनालिका} $I = \qty{1.5}{A}$ ले जाती है।
$n$ निकालिए, फिर भीतर $B$; धारा दुगुनी करने पर $B$ क्या हो जाता
है?
\end{exercise}

\begin{exercise}[$\star$]\label{exo:g11:magnetic-fields:5}
लंबाई \qty{10}{cm} का तार-खंड \qty{0.20}{T} के क्षेत्र के लंबवत \qty{5.0}{A} ले
जाता है। \omterm{prop:g11:magnetic-fields:laplace}{लाप्लास बल} निकालिए; वह कितने द्रव्यमान के भार के
बराबर है?
\end{exercise}

\begin{exercise}[$\star\star$]\label{exo:g11:magnetic-fields:6}
जहाँ एक पदयात्री खड़ी है वहाँ \omterm{def:g11:magnetic-fields:earth}{दिक्पात} $\ang{8}$ पश्चिम है। वह
``कुतुबनुमा के उत्तर'' की ओर \qty{2.0}{km} चलती है: वह सच्चे उत्तर से किस ओर और
कितना बह जाएगी? किस दिशा पर चलने से वह सच्चे उत्तर पहुँचती?
\end{exercise}

\begin{exercise}[$\star\star$]\label{exo:g11:magnetic-fields:7}
\qty{20}{A} ले जाते सीधे तार से कितनी दूरी पर उसका क्षेत्र पृथ्वी के
\qty{50}{\micro T} के बराबर हो जाता है? इससे जीवित केबलों के पास कुतुबनुमा पर भरोसा
करने के बारे में क्या निकलता है?
\end{exercise}

\begin{exercise}[$\star\star$]\label{exo:g11:magnetic-fields:8}
$I = \qty{4.0}{A}$ के साथ $B = \qty{10}{mT}$ देती \omterm{def:g11:magnetic-fields:solenoid}{परिनालिका} डिज़ाइन कीजिए: पहले ज़रूरी
$n$ निकालिए, फिर \qty{25}{cm} की \omterm{def:g11:magnetic-fields:solenoid}{कुंडली} के लिए फेरों की संख्या।
\end{exercise}

\begin{exercise}[$\star\star$]\label{exo:g11:magnetic-fields:9}
पटरी वाले प्रयोग में छड़ (द्रव्यमान \qty{20}{g}, पटरियों के बीच लंबाई $L = \qty{12}{cm}$)
\qty{0.50}{T} के लंबवत क्षेत्र में \qty{10}{A} ले जाती है।
\omterm{prop:g11:magnetic-fields:laplace}{लाप्लास बल} निकालिए, फिर छड़ का आरंभिक त्वरण; $g$ से तुलना कीजिए।
\end{exercise}

\begin{exercise}[$\star\star$]\label{exo:g11:magnetic-fields:10}
द्रव्यमान \qty{5.0}{g} और लंबाई \qty{20}{cm} का क्षैतिज तार अपने लंबवत क्षैतिज
क्षेत्र $B = \qty{0.10}{T}$ में रखा है। कितनी धारा पर \omterm{prop:g11:magnetic-fields:laplace}{लाप्लास बल} तार
के भार को काट देगा, यानी वह हवा में ठहर जाएगा? बल किस ओर
होना चाहिए?
\end{exercise}

\begin{exercise}[$\star\star$]\label{exo:g11:magnetic-fields:11}
\omterm{prop:g11:magnetic-fields:laplace}{लाप्लास बल} की दिशा बताइए (या लिखिए कि वह लुप्त हो जाता है):
(क) धारा पन्ने में दाईं ओर, $\vect B$ पन्ने के भीतर; (ख) धारा पन्ने के ऊपर की
ओर, $\vect B$ पन्ने से बाहर; (ग) धारा $\vect B$ के समांतर। चपटा दायाँ हाथ
इस्तेमाल कीजिए।
\end{exercise}

\begin{exercise}[$\star\star\star$]\label{exo:g11:magnetic-fields:12}
एक क्षैतिज तार भौगोलिक उत्तर--दक्षिण दिशा में, कुतुबनुमा से \qty{3.0}{cm} ऊपर, बिछा
है; वहाँ पृथ्वी का क्षैतिज घटक \qty{20}{\micro T} है। तार $I = \qty{5.0}{A}$ ले जाता है:
कुतुबनुमा पर उसका क्षेत्र निकालिए, उसकी दिशा बताइए, फिर सुई का
विक्षेप (क्षेत्र सदिशों की तरह जुड़ते हैं)।
\end{exercise}

\begin{exercise}[$\star\star\star$]\label{exo:g11:magnetic-fields:13}
कबाड़खाने का \omterm{def:g11:magnetic-fields:electromagnet}{विद्युत्चुंबक} \qty{20}{cm} पर लपेटे $400$ फेरों की
\omterm{def:g11:magnetic-fields:solenoid}{कुंडली} है, प्रतिरोध \qty{2.0}{\ohm}, और \qty{12}{V} की आपूर्ति पर; उसका
लोहे का क्रोड नंगे क्षेत्र को $100$ गुना कर देता है। धारा,
नंगी \omterm{def:g11:magnetic-fields:solenoid}{कुंडली} का क्षेत्र, \omterm{def:g10:refraction:fiber}{क्रोड} सहित क्षेत्र और
\omterm{def:g11:magnetic-fields:solenoid}{कुंडली} द्वारा क्षय की गई शक्ति निकालिए (\cref{ch:g11:circuits-and-power})। स्विच
खुलते ही कबाड़ उसी क्षण क्यों गिर जाता है?
\end{exercise}

\begin{exercise}[$\star\star\star$]\label{exo:g11:magnetic-fields:14}
किसी मोटर के आयताकार फंदे में $N = 50$ फेरे हैं जो $I = \qty{2.0}{A}$ ले जाते हैं; उसकी
\qty{5.0}{cm} लंबाई वाली दोनों भुजाएँ क्षेत्र $B = \qty{0.30}{T}$ के लंबवत हैं। इनमें
से हर भुजा पर लगता बल निकालिए; दोनों बल विपरीत क्यों हैं,
और वे फंदे का क्या करते हैं? बाक़ी दोनों भुजाओं पर कभी-कभी कोई बल
क्यों नहीं लगता? हर आधे चक्कर पर धारा उलटनी क्यों पड़ती है?
\end{exercise}

\begin{exercise}[$\star\star\star$]\label{exo:g11:magnetic-fields:15}
\qty{100}{A} ले जाती केबल के कितने पास जाना पड़ेगा कि उसका क्षेत्र
\qty{1}{T} तक पहुँच जाए? केबल की मिलीमीटर भर की त्रिज्या से तुलना कीजिए और
निष्कर्ष निकालिए कि \omterm{def:g11:magnetic-fields:bfield}{टेस्ला} जितने प्रबल क्षेत्र अकेले तारों से नहीं,
परिनालिकाओं से क्यों बनाए जाते हैं। फिर: न्यूट्रॉन तारे की सतह का
\qty{e8}{T} वाला क्षेत्र पृथ्वी से दस की कितनी घातें ऊपर है?
\end{exercise}

\section{समस्या: कुतुबनुमा, क्रेन और MRI}

\begin{problem}[{सप्ताहांत समस्या --- काम पर लगे तीन चुंबक: महासागर पार करता
एक कुतुबनुमा, कार उठाता एक विद्युत्चुंबक, और MRI की एक परिनालिका --- एक ही
क्षेत्र, दस की दस घातों में फैला, तीन काम करता हुआ}]
\label{pb:g11:magnetic-fields:1}
वही सदिश $\vect B$ \qty{50}{\micro T} से पालवाली नाव की दिशा तय करता है, \qty{1}{T} से
कबाड़ लोहा उठाता है और \qty{1.5}{T} से घुटने का चित्र बनाता है --- इस अध्याय के
तीन सूत्रों के लिए तीन मशीनें।

\medskip
\noindent\textbf{भाग I --- कुतुबनुमा।}
जहाज़ की जगह पर पृथ्वी के क्षेत्र का क्षैतिज घटक \qty{20}{\micro T} और
\omterm{def:g10:universal-gravitation:weight}{ऊर्ध्वाधर} घटक \qty{44}{\micro T} है; \omterm{def:g11:magnetic-fields:earth}{दिक्पात} $\ang{6}$ पश्चिम है।
\begin{enumerate}
  \item कुतुबनुमा की सुई (चुंबकीय) उत्तर की ओर संकेत करती ही क्यों है?
  \item कुल क्षेत्र का परिमाण और क्षैतिज से उसका कोण निकालिए।
  \item ``कुतुबनुमा का उत्तर'' पकड़कर \qty{30}{km} चलने पर जहाज़ सच्चे उत्तर वाले
        रास्ते से कितना पश्चिम जा पहुँचेगा?
  \item इस्पात का माल सुई को $\ang{4}$ और पश्चिम खींच देता है: वही प्रश्न।
  \item भूचुंबकीय ध्रुव के पास कुतुबनुमा बेकार हो जाता है: इसका दोषी कौन-सा
        घटक है, और क्यों?
\end{enumerate}

\medskip
\noindent\textbf{भाग II --- क्रेन।}
कबाड़खाने का \omterm{def:g11:magnetic-fields:electromagnet}{विद्युत्चुंबक} \qty{30}{cm} पर लपेटे $600$ फेरों की
\omterm{def:g11:magnetic-fields:solenoid}{परिनालिका} है, प्रतिरोध \qty{1.5}{\ohm}, और \qty{12}{V} की आपूर्ति से
खिलाई जाती है; उसका लोहे का क्रोड नंगे क्षेत्र को $50$ गुना कर
देता है।
\begin{enumerate}[resume]
  \item फेरा-घनत्व $n$ निकालिए।
  \item \omterm{def:g11:magnetic-fields:solenoid}{कुंडली} में धारा निकालिए (\cref{ch:g11:circuits-and-power})।
  \item नंगी \omterm{def:g11:magnetic-fields:solenoid}{कुंडली} का क्षेत्र $\mu_0 n I$ निकालिए।
  \item \omterm{def:g10:refraction:fiber}{क्रोड} सहित क्षेत्र निकालिए; नियोडिमियम के \qty{0.5}{T} से तुलना
        कीजिए।
  \item \omterm{def:g11:magnetic-fields:solenoid}{कुंडली} में क्षय हुई शक्ति निकालिए, फिर \qty{10}{\minute} की
        एक उठाई-पाली की ऊर्जा।
  \item इतना प्रबल स्थायी \omterm{def:g11:magnetic-fields:magnet}{चुंबक} मौजूद है। फिर भी
        \omterm{def:g11:magnetic-fields:electromagnet}{विद्युत्चुंबक} ही क्यों?
\end{enumerate}

\medskip
\noindent\textbf{भाग III --- चरखी की मोटर।}
क्रेन की चरखी दिष्ट धारा मोटर है: $N = 100$ फेरों का आयताकार फंदा, जिसकी
\qty{4.0}{cm} लंबी भुजाएँ $B = \qty{0.25}{T}$ के लंबवत हैं और जो $I = \qty{3.0}{A}$ ले जाता है।
\begin{enumerate}[resume]
  \item $\vect B$ के लंबवत हर भुजा-गुच्छे पर लगता \omterm{prop:g11:magnetic-fields:laplace}{लाप्लास बल}
        निकालिए।
  \item दोनों बल बराबर और विपरीत हैं, फिर भी उनका असर नहीं कटता। वे
        क्या पैदा करते हैं, और उनका \emph{विपरीत} होना ज़रूरी क्यों है?
  \item बाक़ी दोनों भुजाओं पर, जब वे $\vect B$ के समांतर हों, कौन-सा
        बल लगता है?
  \item आधे चक्कर बाद वही बल घुमाव को उलट देंगे। \omterm{rem:g11:magnetic-fields:motor}{दिक्परिवर्तक}
        क्या करता है, और कब?
  \item डिज़ाइन में तीन अलग-अलग ऐसे बदलाव बताइए जिनमें से हर एक बल
        के युग्म को दुगुना कर दे।
\end{enumerate}

\medskip
\noindent\textbf{भाग IV --- MRI।}
यंत्र की \omterm{def:g11:magnetic-fields:solenoid}{परिनालिका} $I = \qty{500}{A}$ की धारा से \omterm{def:g11:electric-gravitational-fields:capacitor}{एकसमान} \qty{1.5}{T} बनाए रखती है।
\begin{enumerate}[resume]
  \item \qty{1.5}{T} पृथ्वी के क्षेत्र का कितना गुना है?
  \item लोहे के किसी \omterm{def:g10:refraction:fiber}{क्रोड} के \emph{बिना} ज़रूरी फेरा-घनत्व $n$ निकालिए।
  \item यदि लपेट का प्रतिरोध \qty{0.20}{\ohm} होता तो वह कितनी
        शक्ति क्षय करती? असली लपेट का कौन-सा गुण इससे बचा लेता है,
        और किस अत्यधिक ठंड की क़ीमत पर?
  \item और अंतिम चोट: जहाज़, क्रेन और MRI के क्षेत्रों को
        अंतरतारकीय अवकाश (\qty{e-10}{T}) से न्यूट्रॉन तारे (\qty{e8}{T}) तक फैली दस की
        घातों वाली सीढ़ी पर रखिए; इन मशीनों के बीच बदला क्या --- भौतिकी, या
        \omterm{def:g11:circuits-and-power:current}{ऐंपियर}?
\end{enumerate}
\end{problem}
